\section{Related Work}
\label{sec:related}

\noindent\textbf{P\&ID digitization.}
The dominant paradigm for P\&ID recognition has long been to chain
specialist modules: a CNN detects and classifies symbols~\cite{rahul2019,
mani2020,cha2019}, a text recognition stage extracts instrument
tags~\cite{baek2019craft}, and a Hough-based line detector reconstructs
pipe connections~\cite{stuermer2023}.
These pipelines are fragile in practice because errors accumulate
across stages and the final graph is assembled without any global
spatial reasoning.
Jamieson~\etal~\cite{jamieson2024} survey the field and point to
data scarcity as the central roadblock, a view echoed across the literature.
The most relevant prior work for us is Stürmer~\etal~\cite{sturmer2025},
who apply Relationformer end-to-end to the OPEN100 benchmark and gain
more than 25 percentage points in edge detection over the modular
baseline --- but require 60 annotated real P\&IDs to do so.
Our work asks what is possible when that annotated corpus is unavailable.

\noindent\textbf{Synthetic data for engineering diagrams.}
Paliwal~\etal~\cite{paliwal2021} released Dataset-P\&ID, a collection
of 500 synthetic diagrams built by randomly scattering 32 symbol templates
across a blank canvas.
Stürmer~\etal~\cite{sturmer2025} use a similar strategy for pre-training
their model, but conclude in Appendix~A5 that synth-only training
``suffers significantly'' on real images.
GAN-based augmentation~\cite{nurminen2020,elyan2020} can diversify the
visual appearance of symbols but cannot alter graph connectivity,
so it does not address the structural mismatch we identify.
The core issue with random placement is that the resulting graphs look
nothing like real process plants in terms of degree distribution,
component clustering, or flow directionality.
Our generator is specifically designed to avoid this by inheriting
topology from real seeds.

\noindent\textbf{Synthetic-to-real transfer.}
Domain transfer from simulation to real data has a long history in
computer vision~\cite{tobin2017,ganin2016}, and lessons from that
literature carry over here.
Work on robotic manipulation, for instance, has shown that matching the
structural properties of the simulated environment --- physical dynamics,
joint constraints --- matters far more than photographic
realism~\cite{akkaya2019}.
The analogy to P\&IDs is direct: graph topology plays the role of
physical dynamics, and getting it right is what enables transfer.

\noindent\textbf{Image-to-graph prediction.}
Relationformer~\cite{shit2022} builds on Deformable DETR~\cite{zhu2021}
by introducing a shared relation token that, combined pairwise with
object queries, produces joint detections and edge predictions.
More recently, EGTR~\cite{im2024} avoids explicit relation tokens by
reading scene-graph structure from DETR's own self-attention maps,
achieving strong results on Visual Genome.
We use Relationformer to keep our architecture identical to that
of~\cite{sturmer2025}, so that any performance differences can be
attributed cleanly to training data rather than modelling choices.
